\chapter{V2 : rayon de bord d'attaque libre et XFOIL de référence}\label{ch:v2}
\iterbanner{V2}

\begin{evolution}{V2}{pourquoi ce chapitre}La V1 a montré deux limites du modèle de la note d'intention. Avec la famille NACA à 4 chiffres, l'épaisseur pilote à la fois la traînée et le rayon de bord d'attaque, donc la SER : il n'y a pas de vrai compromis. Et notre modèle rapide surestime $C_D$ de 27~\% par rapport à XFOIL, au point de déplacer l'optimum. La V2 change donc la géométrie et le simulateur de référence. Les chapitres~\ref{ch:intro} à \ref{ch:validation} restent valables pour le modèle rapide, conservé comme outil de développement.\end{evolution}

\section{La série NACA à 4 chiffres modifiée}
La série modifiée (Abbott et von Doenhoff, §6.4) garde la ligne moyenne à deux paraboles du chapitre~\ref{ch:naca}, mais remplace la loi d'épaisseur par deux polynômes raccordés à l'abscisse $T$ de l'épaisseur maximale, avec un nouveau paramètre : l'indice de rayon de bord d'attaque $I$. Pour un profil de 20~\% (demi-épaisseur 0,1 en $x=T$) :
\[
y_{20}(x)=\begin{cases}a_0\sqrt x+a_1x+a_2x^2+a_3x^3 & x\le T,\\ d_1(1-x)+d_2(1-x)^2+d_3(1-x)^3 & x\ge T,\end{cases}
\qquad y_t=\frac{t}{0{,}2}\,y_{20}.
\]
\begin{itemize}
\item $a_0=0{,}2969\,I/6$ fixe le nez. Le rayon de bord d'attaque vaut $r_{BA}=1{,}1019\,(tI/6)^2$ : à épaisseur donnée, il varie comme $I^2$. $I=6$ redonne le rayon du NACA standard.
\item $d_1=\dfrac{2{,}24-5{,}42T+12{,}3T^2}{10(1-0{,}878T)}$ fixe la pente au bord de fuite (formule de Ladson et Brooks, 1975), et $d_0=0$ ferme le bord de fuite.
\item $d_2$, $d_3$ sont choisis pour que $y_{20}(T)=0{,}1$ et $y_{20}'(T)=0$ ; puis $a_1$, $a_2$, $a_3$ pour que l'arc avant ait aussi $y=0{,}1$, $y'=0$ \emph{et la même courbure} que l'arc arrière en $x=T$. Deux petits systèmes linéaires suffisent.
\end{itemize}
Avec $I=6$ et $T=0{,}3$, on retrouve la loi du NACA standard à 0,1~\% de corde près. En V2, $T=0{,}3$ est fixé, et $I$ devient la quatrième variable de conception.

\begin{figure}[htbp]
\centering
% figure générée par tools/make_figures_v2.py : v2_formes
\begin{tikzpicture}
\begin{groupplot}[group style={group size=2 by 1,horizontal sep=1.4cm},title style={font=\small},tick label style={font=\scriptsize}]
\nextgroupplot[width=0.62\linewidth,height=4.2cm,axis equal image,xmin=-0.02,xmax=1.02,ymin=-0.08,ymax=0.12,
  xlabel={$x/c$},ylabel={$y/c$},title={NACA 2412 modifié, quatre rayons de bord d'attaque},
  legend style={draw=none,font=\scriptsize,at={(0.5,-0.45)},anchor=north},legend columns=2]
\addplot[corange,very thick] coordinates {(1,0) (0.99989,2.2648e-05) (0.99957,9.0569e-05) (0.99902,0.0002037) (0.99826,0.00036191) (0.99729,0.00056507) (0.99609,0.00081295) (0.99469,0.0011053) (0.99306,0.0014419) (0.99122,0.0018223) (0.98917,0.0022462) (0.98691,0.0027131) (0.98443,0.0032226) (0.98175,0.0037741) (0.97885,0.0043671) (0.97575,0.005001) (0.97243,0.0056751) (0.96892,0.0063887) (0.96519,0.0071412) (0.96127,0.0079316) (0.95714,0.0087592) (0.95282,0.0096231) (0.94829,0.010522) (0.94358,0.011456) (0.93866,0.012423) (0.93356,0.013423) (0.92826,0.014454) (0.92278,0.015515) (0.91712,0.016605) (0.91127,0.017723) (0.90524,0.018867) (0.89903,0.020037) (0.89265,0.021231) (0.88609,0.022448) (0.87937,0.023685) (0.87248,0.024943) (0.86542,0.026219) (0.85821,0.027512) (0.85083,0.02882) (0.84331,0.030143) (0.83563,0.031477) (0.8278,0.032823) (0.81982,0.034178) (0.81171,0.03554) (0.80346,0.036908) (0.79507,0.03828) (0.78655,0.039655) (0.77791,0.041031) (0.76914,0.042406) (0.76025,0.043778) (0.75124,0.045146) (0.74213,0.046509) (0.7329,0.047864) (0.72357,0.049209) (0.71414,0.050544) (0.70461,0.051865) (0.695,0.053173) (0.68529,0.054465) (0.6755,0.055739) (0.66563,0.056993) (0.65569,0.058227) (0.64568,0.059439) (0.6356,0.060626) (0.62546,0.061789) (0.61526,0.062924) (0.60501,0.064031) (0.59471,0.065108) (0.58436,0.066155) (0.57398,0.067169) (0.56356,0.06815) (0.55312,0.069096) (0.54264,0.070006) (0.53215,0.070879) (0.52164,0.071714) (0.51112,0.072511) (0.5006,0.073268) (0.49007,0.073984) (0.47955,0.074659) (0.46903,0.075293) (0.45852,0.075884) (0.44804,0.076431) (0.43757,0.076936) (0.42713,0.077397) (0.41672,0.077813) (0.40635,0.078186) (0.39599,0.078513) (0.38562,0.078784) (0.3753,0.078996) (0.36503,0.07915) (0.35482,0.079245) (0.34468,0.079283) (0.3346,0.079265) (0.3246,0.07919) (0.31468,0.079061) (0.30484,0.078878) (0.29508,0.078642) (0.28542,0.078353) (0.27585,0.078004) (0.26639,0.077593) (0.25702,0.077115) (0.24777,0.076566) (0.23864,0.075944) (0.22962,0.075246) (0.22073,0.074469) (0.21196,0.073614) (0.20333,0.072678) (0.19483,0.071661) (0.18647,0.070562) (0.17826,0.069384) (0.1702,0.068125) (0.16229,0.066788) (0.15453,0.065375) (0.14694,0.063887) (0.13951,0.062328) (0.13225,0.060701) (0.12515,0.059009) (0.11823,0.057255) (0.11149,0.055446) (0.10493,0.053583) (0.098544,0.051674) (0.092346,0.049722) (0.086336,0.047734) (0.080515,0.045714) (0.074885,0.043669) (0.069449,0.041604) (0.064208,0.039526) (0.059164,0.03744) (0.054319,0.035353) (0.049675,0.033271) (0.045232,0.0312) (0.040992,0.029148) (0.036957,0.027118) (0.033127,0.025119) (0.029503,0.023156) (0.026086,0.021234) (0.022877,0.019359) (0.019877,0.017538) (0.017086,0.015775) (0.014504,0.014074) (0.012133,0.012442) (0.0099725,0.010882) (0.008023,0.0093986) (0.0062846,0.0079955) (0.0047577,0.0066763) (0.0034425,0.0054439) (0.002339,0.0043012) (0.0014474,0.0032504) (0.00076767,0.0022936) (0.00029989,0.0014323) (4.4022e-05,0.00066751) (0,0) (0.00017529,-0.00064558) (0.00057728,-0.0013446) (0.0012056,-0.0020965) (0.0020598,-0.0029005) (0.0031391,-0.0037552) (0.0044428,-0.0046592) (0.0059699,-0.0056107) (0.0077193,-0.0066074) (0.0096898,-0.007647) (0.01188,-0.0087267) (0.014288,-0.0098435) (0.016913,-0.010994) (0.019752,-0.012176) (0.022804,-0.013384) (0.026067,-0.014615) (0.029538,-0.015865) (0.033215,-0.01713) (0.037097,-0.018405) (0.04118,-0.019686) (0.045462,-0.020969) (0.049941,-0.022249) (0.054614,-0.023522) (0.059477,-0.024783) (0.064529,-0.026027) (0.069767,-0.02725) (0.075187,-0.028448) (0.080787,-0.029616) (0.086564,-0.030751) (0.092515,-0.031848) (0.098637,-0.032904) (0.10493,-0.033915) (0.11138,-0.034877) (0.118,-0.035788) (0.12477,-0.036646) (0.1317,-0.037446) (0.13879,-0.038189) (0.14602,-0.038871) (0.1534,-0.039491) (0.16092,-0.040049) (0.16858,-0.040543) (0.17638,-0.040974) (0.18432,-0.041342) (0.19238,-0.041648) (0.20057,-0.041892) (0.20889,-0.042076) (0.21733,-0.042203) (0.22588,-0.042275) (0.23455,-0.042294) (0.24334,-0.042265) (0.25223,-0.042191) (0.26122,-0.042076) (0.27032,-0.041926) (0.27951,-0.041745) (0.2888,-0.041539) (0.29818,-0.041314) (0.30765,-0.041076) (0.3172,-0.040828) (0.32683,-0.04057) (0.33653,-0.040303) (0.3463,-0.040026) (0.35615,-0.039741) (0.36605,-0.039446) (0.37601,-0.039144) (0.38603,-0.038834) (0.3961,-0.038517) (0.40627,-0.03819) (0.41651,-0.037844) (0.42679,-0.037477) (0.4371,-0.037091) (0.44744,-0.036685) (0.4578,-0.036259) (0.46818,-0.035816) (0.47858,-0.035354) (0.48899,-0.034875) (0.4994,-0.034379) (0.50982,-0.033867) (0.52023,-0.03334) (0.53064,-0.032797) (0.54103,-0.032241) (0.55141,-0.031672) (0.56177,-0.03109) (0.5721,-0.030496) (0.58241,-0.029892) (0.59268,-0.029277) (0.60291,-0.028653) (0.61309,-0.028021) (0.62323,-0.027381) (0.63332,-0.026734) (0.64335,-0.026082) (0.65333,-0.025424) (0.66323,-0.024763) (0.67307,-0.024098) (0.68283,-0.02343) (0.69252,-0.022761) (0.70212,-0.022091) (0.71164,-0.021421) (0.72106,-0.020752) (0.7304,-0.020085) (0.73963,-0.01942) (0.74876,-0.018757) (0.75778,-0.018099) (0.76669,-0.017446) (0.77549,-0.016797) (0.78416,-0.016155) (0.79272,-0.015519) (0.80114,-0.01489) (0.80944,-0.01427) (0.8176,-0.013658) (0.82562,-0.013054) (0.8335,-0.01246) (0.84124,-0.011877) (0.84883,-0.011304) (0.85626,-0.010742) (0.86354,-0.010191) (0.87067,-0.009652) (0.87763,-0.0091254) (0.88442,-0.0086114) (0.89104,-0.0081104) (0.8975,-0.0076226) (0.90378,-0.0071484) (0.90988,-0.0066881) (0.91581,-0.0062418) (0.92155,-0.0058099) (0.9271,-0.0053925) (0.93247,-0.0049899) (0.93764,-0.0046021) (0.94263,-0.0042295) (0.94742,-0.003872) (0.95201,-0.0035299) (0.9564,-0.0032032) (0.96059,-0.0028921) (0.96458,-0.0025966) (0.96837,-0.0023168) (0.97194,-0.0020528) (0.97531,-0.0018046) (0.97847,-0.0015723) (0.98142,-0.0013559) (0.98415,-0.0011555) (0.98667,-0.00097106) (0.98897,-0.0008026) (0.99106,-0.00065015) (0.99293,-0.00051373) (0.99459,-0.00039334) (0.99602,-0.00028899) (0.99724,-0.00020069) (0.99823,-0.00012845) (0.999,-7.2252e-05) (0.99956,-3.2112e-05) (0.99989,-8.0281e-06) (1,0)};\addlegendentry{$I=2$, $r_{BA}=0.18$~\%}
\addplot[cgreen,very thick] coordinates {(1,0) (0.99989,2.2648e-05) (0.99957,9.0569e-05) (0.99902,0.0002037) (0.99826,0.00036191) (0.99729,0.00056507) (0.99609,0.00081295) (0.99469,0.0011053) (0.99306,0.0014419) (0.99122,0.0018223) (0.98917,0.0022462) (0.98691,0.0027131) (0.98443,0.0032226) (0.98175,0.0037741) (0.97885,0.0043671) (0.97575,0.005001) (0.97243,0.0056751) (0.96892,0.0063887) (0.96519,0.0071412) (0.96127,0.0079316) (0.95714,0.0087592) (0.95282,0.0096231) (0.94829,0.010522) (0.94358,0.011456) (0.93866,0.012423) (0.93356,0.013423) (0.92826,0.014454) (0.92278,0.015515) (0.91712,0.016605) (0.91127,0.017723) (0.90524,0.018867) (0.89903,0.020037) (0.89265,0.021231) (0.88609,0.022448) (0.87937,0.023685) (0.87248,0.024943) (0.86542,0.026219) (0.85821,0.027512) (0.85083,0.02882) (0.84331,0.030143) (0.83563,0.031477) (0.8278,0.032823) (0.81982,0.034178) (0.81171,0.03554) (0.80346,0.036908) (0.79507,0.03828) (0.78655,0.039655) (0.77791,0.041031) (0.76914,0.042406) (0.76025,0.043778) (0.75124,0.045146) (0.74213,0.046509) (0.7329,0.047864) (0.72357,0.049209) (0.71414,0.050544) (0.70461,0.051865) (0.695,0.053173) (0.68529,0.054465) (0.6755,0.055739) (0.66563,0.056993) (0.65569,0.058227) (0.64568,0.059439) (0.6356,0.060626) (0.62546,0.061789) (0.61526,0.062924) (0.60501,0.064031) (0.59471,0.065108) (0.58436,0.066155) (0.57398,0.067169) (0.56356,0.06815) (0.55312,0.069096) (0.54264,0.070006) (0.53215,0.070879) (0.52164,0.071714) (0.51112,0.072511) (0.5006,0.073268) (0.49007,0.073984) (0.47955,0.074659) (0.46903,0.075293) (0.45852,0.075884) (0.44804,0.076431) (0.43757,0.076936) (0.42713,0.077397) (0.41672,0.077813) (0.40635,0.078186) (0.39599,0.078513) (0.38562,0.078784) (0.3753,0.078996) (0.36503,0.07915) (0.35482,0.079245) (0.34468,0.079283) (0.3346,0.079265) (0.3246,0.07919) (0.31468,0.079061) (0.30484,0.078878) (0.29508,0.078642) (0.28542,0.078353) (0.27585,0.078008) (0.26638,0.077605) (0.25702,0.07714) (0.24777,0.076612) (0.23863,0.076019) (0.22961,0.075359) (0.22072,0.074633) (0.21195,0.073837) (0.20331,0.072973) (0.19481,0.072041) (0.18645,0.071039) (0.17823,0.069969) (0.17016,0.068832) (0.16224,0.067628) (0.15447,0.066359) (0.14687,0.065027) (0.13943,0.063634) (0.13215,0.062181) (0.12504,0.060672) (0.1181,0.059109) (0.11134,0.057495) (0.10476,0.055833) (0.098361,0.054127) (0.092144,0.052379) (0.086114,0.050594) (0.080272,0.048775) (0.074622,0.046927) (0.069166,0.045052) (0.063905,0.043156) (0.058842,0.041243) (0.053979,0.039316) (0.049317,0.03738) (0.044859,0.035439) (0.040604,0.033497) (0.036556,0.031558) (0.032715,0.029627) (0.029083,0.027707) (0.025661,0.025803) (0.022449,0.023918) (0.019449,0.022056) (0.016662,0.02022) (0.014087,0.018413) (0.011727,0.01664) (0.0095815,0.014902) (0.0076509,0.013204) (0.0059358,0.011546) (0.0044365,0.009932) (0.0031534,0.0083637) (0.0020865,0.006843) (0.0012362,0.0053715) (0.00060239,0.0039505) (0.00018514,0.0025811) (-1.5618e-05,0.0012641) (0,0) (0.00023493,-0.0012421) (0.00069203,-0.0024934) (0.0013709,-0.0037534) (0.0022709,-0.0050215) (0.0033916,-0.0062971) (0.0047319,-0.0075791) (0.0062911,-0.0088664) (0.0080682,-0.010158) (0.010062,-0.011452) (0.012271,-0.012747) (0.014694,-0.014041) (0.017329,-0.015333) (0.020176,-0.016621) (0.023232,-0.017901) (0.026494,-0.019173) (0.029963,-0.020434) (0.033635,-0.021681) (0.037508,-0.022913) (0.041581,-0.024126) (0.04585,-0.025318) (0.050314,-0.026487) (0.054971,-0.027631) (0.059817,-0.028745) (0.064851,-0.02983) (0.070069,-0.030881) (0.07547,-0.031896) (0.08105,-0.032874) (0.086806,-0.033812) (0.092737,-0.034708) (0.098839,-0.03556) (0.10511,-0.036367) (0.11154,-0.037127) (0.11814,-0.037838) (0.1249,-0.0385) (0.13181,-0.03911) (0.13888,-0.039669) (0.1461,-0.040177) (0.15347,-0.040631) (0.16098,-0.041033) (0.16863,-0.041383) (0.17642,-0.041681) (0.18435,-0.041928) (0.1924,-0.042124) (0.20059,-0.042272) (0.2089,-0.042372) (0.21734,-0.042427) (0.22589,-0.042438) (0.23456,-0.042408) (0.24334,-0.04234) (0.25223,-0.042237) (0.26122,-0.042101) (0.27032,-0.041937) (0.27951,-0.041749) (0.2888,-0.041539) (0.29818,-0.041314) (0.30765,-0.041076) (0.3172,-0.040828) (0.32683,-0.04057) (0.33653,-0.040303) (0.3463,-0.040026) (0.35615,-0.039741) (0.36605,-0.039446) (0.37601,-0.039144) (0.38603,-0.038834) (0.3961,-0.038517) (0.40627,-0.03819) (0.41651,-0.037844) (0.42679,-0.037477) (0.4371,-0.037091) (0.44744,-0.036685) (0.4578,-0.036259) (0.46818,-0.035816) (0.47858,-0.035354) (0.48899,-0.034875) (0.4994,-0.034379) (0.50982,-0.033867) (0.52023,-0.03334) (0.53064,-0.032797) (0.54103,-0.032241) (0.55141,-0.031672) (0.56177,-0.03109) (0.5721,-0.030496) (0.58241,-0.029892) (0.59268,-0.029277) (0.60291,-0.028653) (0.61309,-0.028021) (0.62323,-0.027381) (0.63332,-0.026734) (0.64335,-0.026082) (0.65333,-0.025424) (0.66323,-0.024763) (0.67307,-0.024098) (0.68283,-0.02343) (0.69252,-0.022761) (0.70212,-0.022091) (0.71164,-0.021421) (0.72106,-0.020752) (0.7304,-0.020085) (0.73963,-0.01942) (0.74876,-0.018757) (0.75778,-0.018099) (0.76669,-0.017446) (0.77549,-0.016797) (0.78416,-0.016155) (0.79272,-0.015519) (0.80114,-0.01489) (0.80944,-0.01427) (0.8176,-0.013658) (0.82562,-0.013054) (0.8335,-0.01246) (0.84124,-0.011877) (0.84883,-0.011304) (0.85626,-0.010742) (0.86354,-0.010191) (0.87067,-0.009652) (0.87763,-0.0091254) (0.88442,-0.0086114) (0.89104,-0.0081104) (0.8975,-0.0076226) (0.90378,-0.0071484) (0.90988,-0.0066881) (0.91581,-0.0062418) (0.92155,-0.0058099) (0.9271,-0.0053925) (0.93247,-0.0049899) (0.93764,-0.0046021) (0.94263,-0.0042295) (0.94742,-0.003872) (0.95201,-0.0035299) (0.9564,-0.0032032) (0.96059,-0.0028921) (0.96458,-0.0025966) (0.96837,-0.0023168) (0.97194,-0.0020528) (0.97531,-0.0018046) (0.97847,-0.0015723) (0.98142,-0.0013559) (0.98415,-0.0011555) (0.98667,-0.00097106) (0.98897,-0.0008026) (0.99106,-0.00065015) (0.99293,-0.00051373) (0.99459,-0.00039334) (0.99602,-0.00028899) (0.99724,-0.00020069) (0.99823,-0.00012845) (0.999,-7.2252e-05) (0.99956,-3.2112e-05) (0.99989,-8.0281e-06) (1,0)};\addlegendentry{$I=4$, $r_{BA}=0.71$~\%}
\addplot[cblue,very thick] coordinates {(1,0) (0.99989,2.2648e-05) (0.99957,9.0569e-05) (0.99902,0.0002037) (0.99826,0.00036191) (0.99729,0.00056507) (0.99609,0.00081295) (0.99469,0.0011053) (0.99306,0.0014419) (0.99122,0.0018223) (0.98917,0.0022462) (0.98691,0.0027131) (0.98443,0.0032226) (0.98175,0.0037741) (0.97885,0.0043671) (0.97575,0.005001) (0.97243,0.0056751) (0.96892,0.0063887) (0.96519,0.0071412) (0.96127,0.0079316) (0.95714,0.0087592) (0.95282,0.0096231) (0.94829,0.010522) (0.94358,0.011456) (0.93866,0.012423) (0.93356,0.013423) (0.92826,0.014454) (0.92278,0.015515) (0.91712,0.016605) (0.91127,0.017723) (0.90524,0.018867) (0.89903,0.020037) (0.89265,0.021231) (0.88609,0.022448) (0.87937,0.023685) (0.87248,0.024943) (0.86542,0.026219) (0.85821,0.027512) (0.85083,0.02882) (0.84331,0.030143) (0.83563,0.031477) (0.8278,0.032823) (0.81982,0.034178) (0.81171,0.03554) (0.80346,0.036908) (0.79507,0.03828) (0.78655,0.039655) (0.77791,0.041031) (0.76914,0.042406) (0.76025,0.043778) (0.75124,0.045146) (0.74213,0.046509) (0.7329,0.047864) (0.72357,0.049209) (0.71414,0.050544) (0.70461,0.051865) (0.695,0.053173) (0.68529,0.054465) (0.6755,0.055739) (0.66563,0.056993) (0.65569,0.058227) (0.64568,0.059439) (0.6356,0.060626) (0.62546,0.061789) (0.61526,0.062924) (0.60501,0.064031) (0.59471,0.065108) (0.58436,0.066155) (0.57398,0.067169) (0.56356,0.06815) (0.55312,0.069096) (0.54264,0.070006) (0.53215,0.070879) (0.52164,0.071714) (0.51112,0.072511) (0.5006,0.073268) (0.49007,0.073984) (0.47955,0.074659) (0.46903,0.075293) (0.45852,0.075884) (0.44804,0.076431) (0.43757,0.076936) (0.42713,0.077397) (0.41672,0.077813) (0.40635,0.078186) (0.39599,0.078513) (0.38562,0.078784) (0.3753,0.078996) (0.36503,0.07915) (0.35482,0.079245) (0.34468,0.079283) (0.3346,0.079265) (0.3246,0.07919) (0.31468,0.079061) (0.30484,0.078878) (0.29508,0.078642) (0.28542,0.078354) (0.27585,0.078013) (0.26638,0.077616) (0.25702,0.077165) (0.24777,0.076658) (0.23863,0.076094) (0.22961,0.075473) (0.22071,0.074796) (0.21194,0.074061) (0.2033,0.073269) (0.19479,0.07242) (0.18642,0.071515) (0.1782,0.070554) (0.17012,0.069538) (0.16219,0.068468) (0.15441,0.067344) (0.1468,0.066167) (0.13934,0.06494) (0.13205,0.063662) (0.12493,0.062336) (0.11798,0.060963) (0.1112,0.059545) (0.1046,0.058083) (0.098178,0.056579) (0.091942,0.055036) (0.085892,0.053454) (0.08003,0.051836) (0.07436,0.050184) (0.068884,0.0485) (0.063603,0.046787) (0.058521,0.045046) (0.053639,0.043279) (0.04896,0.041488) (0.044485,0.039677) (0.040216,0.037846) (0.036156,0.035998) (0.032304,0.034135) (0.028664,0.032259) (0.025236,0.030372) (0.022021,0.028476) (0.019021,0.026573) (0.016237,0.024665) (0.01367,0.022752) (0.011321,0.020838) (0.0091904,0.018923) (0.0072788,0.017008) (0.005587,0.015096) (0.0041153,0.013188) (0.0028643,0.011284) (0.0018341,0.0093849) (0.001025,0.0074925) (0.00043712,0.0056073) (7.0386e-05,0.0037299) (-7.5258e-05,0.0018606) (0,0) (0.00029457,-0.0018387) (0.00080678,-0.0036422) (0.0015362,-0.0054103) (0.0024821,-0.0071426) (0.003644,-0.0088389) (0.005021,-0.010499) (0.0066124,-0.012122) (0.008417,-0.013708) (0.010434,-0.015257) (0.012662,-0.016767) (0.0151,-0.018239) (0.017746,-0.019672) (0.0206,-0.021066) (0.023659,-0.022419) (0.026922,-0.023732) (0.030388,-0.025003) (0.034054,-0.026233) (0.037919,-0.027421) (0.041981,-0.028566) (0.046238,-0.029668) (0.050688,-0.030726) (0.055328,-0.031739) (0.060157,-0.032708) (0.065172,-0.033632) (0.070371,-0.034511) (0.075752,-0.035344) (0.081312,-0.036132) (0.087048,-0.036873) (0.092959,-0.037568) (0.099041,-0.038217) (0.10529,-0.03882) (0.11171,-0.039377) (0.11829,-0.039888) (0.12503,-0.040353) (0.13193,-0.040774) (0.13898,-0.04115) (0.14619,-0.041482) (0.15354,-0.041771) (0.16104,-0.042018) (0.16868,-0.042223) (0.17646,-0.042388) (0.18438,-0.042513) (0.19243,-0.042601) (0.20061,-0.042652) (0.20892,-0.042668) (0.21735,-0.04265) (0.2259,-0.042601) (0.23456,-0.042522) (0.24334,-0.042415) (0.25223,-0.042283) (0.26122,-0.042127) (0.27032,-0.041949) (0.27951,-0.041753) (0.2888,-0.04154) (0.29818,-0.041314) (0.30765,-0.041076) (0.3172,-0.040828) (0.32683,-0.04057) (0.33653,-0.040303) (0.3463,-0.040026) (0.35615,-0.039741) (0.36605,-0.039446) (0.37601,-0.039144) (0.38603,-0.038834) (0.3961,-0.038517) (0.40627,-0.03819) (0.41651,-0.037844) (0.42679,-0.037477) (0.4371,-0.037091) (0.44744,-0.036685) (0.4578,-0.036259) (0.46818,-0.035816) (0.47858,-0.035354) (0.48899,-0.034875) (0.4994,-0.034379) (0.50982,-0.033867) (0.52023,-0.03334) (0.53064,-0.032797) (0.54103,-0.032241) (0.55141,-0.031672) (0.56177,-0.03109) (0.5721,-0.030496) (0.58241,-0.029892) (0.59268,-0.029277) (0.60291,-0.028653) (0.61309,-0.028021) (0.62323,-0.027381) (0.63332,-0.026734) (0.64335,-0.026082) (0.65333,-0.025424) (0.66323,-0.024763) (0.67307,-0.024098) (0.68283,-0.02343) (0.69252,-0.022761) (0.70212,-0.022091) (0.71164,-0.021421) (0.72106,-0.020752) (0.7304,-0.020085) (0.73963,-0.01942) (0.74876,-0.018757) (0.75778,-0.018099) (0.76669,-0.017446) (0.77549,-0.016797) (0.78416,-0.016155) (0.79272,-0.015519) (0.80114,-0.01489) (0.80944,-0.01427) (0.8176,-0.013658) (0.82562,-0.013054) (0.8335,-0.01246) (0.84124,-0.011877) (0.84883,-0.011304) (0.85626,-0.010742) (0.86354,-0.010191) (0.87067,-0.009652) (0.87763,-0.0091254) (0.88442,-0.0086114) (0.89104,-0.0081104) (0.8975,-0.0076226) (0.90378,-0.0071484) (0.90988,-0.0066881) (0.91581,-0.0062418) (0.92155,-0.0058099) (0.9271,-0.0053925) (0.93247,-0.0049899) (0.93764,-0.0046021) (0.94263,-0.0042295) (0.94742,-0.003872) (0.95201,-0.0035299) (0.9564,-0.0032032) (0.96059,-0.0028921) (0.96458,-0.0025966) (0.96837,-0.0023168) (0.97194,-0.0020528) (0.97531,-0.0018046) (0.97847,-0.0015723) (0.98142,-0.0013559) (0.98415,-0.0011555) (0.98667,-0.00097106) (0.98897,-0.0008026) (0.99106,-0.00065015) (0.99293,-0.00051373) (0.99459,-0.00039334) (0.99602,-0.00028899) (0.99724,-0.00020069) (0.99823,-0.00012845) (0.999,-7.2252e-05) (0.99956,-3.2112e-05) (0.99989,-8.0281e-06) (1,0)};\addlegendentry{$I=6$, $r_{BA}=1.59$~\%}
\addplot[cmuted,very thick] coordinates {(1,0) (0.99989,2.2648e-05) (0.99957,9.0569e-05) (0.99902,0.0002037) (0.99826,0.00036191) (0.99729,0.00056507) (0.99609,0.00081295) (0.99469,0.0011053) (0.99306,0.0014419) (0.99122,0.0018223) (0.98917,0.0022462) (0.98691,0.0027131) (0.98443,0.0032226) (0.98175,0.0037741) (0.97885,0.0043671) (0.97575,0.005001) (0.97243,0.0056751) (0.96892,0.0063887) (0.96519,0.0071412) (0.96127,0.0079316) (0.95714,0.0087592) (0.95282,0.0096231) (0.94829,0.010522) (0.94358,0.011456) (0.93866,0.012423) (0.93356,0.013423) (0.92826,0.014454) (0.92278,0.015515) (0.91712,0.016605) (0.91127,0.017723) (0.90524,0.018867) (0.89903,0.020037) (0.89265,0.021231) (0.88609,0.022448) (0.87937,0.023685) (0.87248,0.024943) (0.86542,0.026219) (0.85821,0.027512) (0.85083,0.02882) (0.84331,0.030143) (0.83563,0.031477) (0.8278,0.032823) (0.81982,0.034178) (0.81171,0.03554) (0.80346,0.036908) (0.79507,0.03828) (0.78655,0.039655) (0.77791,0.041031) (0.76914,0.042406) (0.76025,0.043778) (0.75124,0.045146) (0.74213,0.046509) (0.7329,0.047864) (0.72357,0.049209) (0.71414,0.050544) (0.70461,0.051865) (0.695,0.053173) (0.68529,0.054465) (0.6755,0.055739) (0.66563,0.056993) (0.65569,0.058227) (0.64568,0.059439) (0.6356,0.060626) (0.62546,0.061789) (0.61526,0.062924) (0.60501,0.064031) (0.59471,0.065108) (0.58436,0.066155) (0.57398,0.067169) (0.56356,0.06815) (0.55312,0.069096) (0.54264,0.070006) (0.53215,0.070879) (0.52164,0.071714) (0.51112,0.072511) (0.5006,0.073268) (0.49007,0.073984) (0.47955,0.074659) (0.46903,0.075293) (0.45852,0.075884) (0.44804,0.076431) (0.43757,0.076936) (0.42713,0.077397) (0.41672,0.077813) (0.40635,0.078186) (0.39599,0.078513) (0.38562,0.078784) (0.3753,0.078996) (0.36503,0.07915) (0.35482,0.079245) (0.34468,0.079283) (0.3346,0.079265) (0.3246,0.07919) (0.31468,0.079061) (0.30484,0.078878) (0.29508,0.078642) (0.28542,0.078355) (0.27585,0.078017) (0.26638,0.077628) (0.25702,0.07719) (0.24777,0.076704) (0.23863,0.076169) (0.22961,0.075587) (0.2207,0.074959) (0.21193,0.074284) (0.20328,0.073564) (0.19477,0.0728) (0.1864,0.071992) (0.17816,0.07114) (0.17008,0.070245) (0.16214,0.069307) (0.15435,0.068328) (0.14673,0.067307) (0.13926,0.066245) (0.13195,0.065143) (0.12481,0.064) (0.11785,0.062817) (0.11105,0.061595) (0.10443,0.060333) (0.097995,0.059032) (0.09174,0.057692) (0.08567,0.056314) (0.079788,0.054897) (0.074098,0.053442) (0.068601,0.051949) (0.063301,0.050418) (0.0582,0.048848) (0.053299,0.047242) (0.048603,0.045597) (0.044112,0.043915) (0.039828,0.042195) (0.035755,0.040438) (0.031893,0.038643) (0.028244,0.036811) (0.02481,0.034942) (0.021593,0.033035) (0.018594,0.031091) (0.015813,0.02911) (0.013254,0.027091) (0.010915,0.025036) (0.0087994,0.022943) (0.0069068,0.020813) (0.0052381,0.018647) (0.0037941,0.016443) (0.0025752,0.014203) (0.0015817,0.011927) (0.00081385,0.0096136) (0.00027184,0.0072642) (-4.4368e-05,0.0048787) (-0.0001349,0.0024572) (0,0) (0.00035421,-0.0024353) (0.00092154,-0.004791) (0.0017014,-0.0070671) (0.0026933,-0.0092636) (0.0038964,-0.011381) (0.0053101,-0.013419) (0.0069336,-0.015378) (0.0087658,-0.017259) (0.010806,-0.019062) (0.013053,-0.020788) (0.015506,-0.022437) (0.018163,-0.024011) (0.021024,-0.025511) (0.024087,-0.026937) (0.02735,-0.02829) (0.030813,-0.029573) (0.034474,-0.030785) (0.038331,-0.031929) (0.042382,-0.033006) (0.046626,-0.034017) (0.051061,-0.034964) (0.055686,-0.035848) (0.060497,-0.036671) (0.065494,-0.037435) (0.070674,-0.038142) (0.076035,-0.038793) (0.081574,-0.039389) (0.087291,-0.039934) (0.093181,-0.040428) (0.099243,-0.040874) (0.10547,-0.041272) (0.11187,-0.041626) (0.11844,-0.041937) (0.12516,-0.042207) (0.13204,-0.042438) (0.13908,-0.042631) (0.14627,-0.042788) (0.15361,-0.042911) (0.1611,-0.043002) (0.16873,-0.043063) (0.1765,-0.043094) (0.18441,-0.043098) (0.19245,-0.043077) (0.20063,-0.043031) (0.20893,-0.042963) (0.21736,-0.042873) (0.2259,-0.042764) (0.23457,-0.042636) (0.24335,-0.04249) (0.25223,-0.042329) (0.26122,-0.042152) (0.27032,-0.041961) (0.27951,-0.041757) (0.2888,-0.041541) (0.29818,-0.041314) (0.30765,-0.041076) (0.3172,-0.040828) (0.32683,-0.04057) (0.33653,-0.040303) (0.3463,-0.040026) (0.35615,-0.039741) (0.36605,-0.039446) (0.37601,-0.039144) (0.38603,-0.038834) (0.3961,-0.038517) (0.40627,-0.03819) (0.41651,-0.037844) (0.42679,-0.037477) (0.4371,-0.037091) (0.44744,-0.036685) (0.4578,-0.036259) (0.46818,-0.035816) (0.47858,-0.035354) (0.48899,-0.034875) (0.4994,-0.034379) (0.50982,-0.033867) (0.52023,-0.03334) (0.53064,-0.032797) (0.54103,-0.032241) (0.55141,-0.031672) (0.56177,-0.03109) (0.5721,-0.030496) (0.58241,-0.029892) (0.59268,-0.029277) (0.60291,-0.028653) (0.61309,-0.028021) (0.62323,-0.027381) (0.63332,-0.026734) (0.64335,-0.026082) (0.65333,-0.025424) (0.66323,-0.024763) (0.67307,-0.024098) (0.68283,-0.02343) (0.69252,-0.022761) (0.70212,-0.022091) (0.71164,-0.021421) (0.72106,-0.020752) (0.7304,-0.020085) (0.73963,-0.01942) (0.74876,-0.018757) (0.75778,-0.018099) (0.76669,-0.017446) (0.77549,-0.016797) (0.78416,-0.016155) (0.79272,-0.015519) (0.80114,-0.01489) (0.80944,-0.01427) (0.8176,-0.013658) (0.82562,-0.013054) (0.8335,-0.01246) (0.84124,-0.011877) (0.84883,-0.011304) (0.85626,-0.010742) (0.86354,-0.010191) (0.87067,-0.009652) (0.87763,-0.0091254) (0.88442,-0.0086114) (0.89104,-0.0081104) (0.8975,-0.0076226) (0.90378,-0.0071484) (0.90988,-0.0066881) (0.91581,-0.0062418) (0.92155,-0.0058099) (0.9271,-0.0053925) (0.93247,-0.0049899) (0.93764,-0.0046021) (0.94263,-0.0042295) (0.94742,-0.003872) (0.95201,-0.0035299) (0.9564,-0.0032032) (0.96059,-0.0028921) (0.96458,-0.0025966) (0.96837,-0.0023168) (0.97194,-0.0020528) (0.97531,-0.0018046) (0.97847,-0.0015723) (0.98142,-0.0013559) (0.98415,-0.0011555) (0.98667,-0.00097106) (0.98897,-0.0008026) (0.99106,-0.00065015) (0.99293,-0.00051373) (0.99459,-0.00039334) (0.99602,-0.00028899) (0.99724,-0.00020069) (0.99823,-0.00012845) (0.999,-7.2252e-05) (0.99956,-3.2112e-05) (0.99989,-8.0281e-06) (1,0)};\addlegendentry{$I=8$, $r_{BA}=2.82$~\%}
\nextgroupplot[width=0.36\linewidth,height=5cm,axis equal image,xmin=-0.002,xmax=0.05,ymin=-0.035,ymax=0.04,
  xlabel={$x/c$},title={zoom sur le nez},xticklabel style={/pgf/number format/fixed}]
\addplot[corange,very thick] coordinates {(0.059164,0.03744) (0.054319,0.035353) (0.049675,0.033271) (0.045232,0.0312) (0.040992,0.029148) (0.036957,0.027118) (0.033127,0.025119) (0.029503,0.023156) (0.026086,0.021234) (0.022877,0.019359) (0.019877,0.017538) (0.017086,0.015775) (0.014504,0.014074) (0.012133,0.012442) (0.0099725,0.010882) (0.008023,0.0093986) (0.0062846,0.0079955) (0.0047577,0.0066763) (0.0034425,0.0054439) (0.002339,0.0043012) (0.0014474,0.0032504) (0.00076767,0.0022936) (0.00029989,0.0014323) (4.4022e-05,0.00066751) (0,0) (0.00017529,-0.00064558) (0.00057728,-0.0013446) (0.0012056,-0.0020965) (0.0020598,-0.0029005) (0.0031391,-0.0037552) (0.0044428,-0.0046592) (0.0059699,-0.0056107) (0.0077193,-0.0066074) (0.0096898,-0.007647) (0.01188,-0.0087267) (0.014288,-0.0098435) (0.016913,-0.010994) (0.019752,-0.012176) (0.022804,-0.013384) (0.026067,-0.014615) (0.029538,-0.015865) (0.033215,-0.01713) (0.037097,-0.018405) (0.04118,-0.019686) (0.045462,-0.020969) (0.049941,-0.022249) (0.054614,-0.023522) (0.059477,-0.024783)};
\addplot[cgreen,very thick] coordinates {(0.058842,0.041243) (0.053979,0.039316) (0.049317,0.03738) (0.044859,0.035439) (0.040604,0.033497) (0.036556,0.031558) (0.032715,0.029627) (0.029083,0.027707) (0.025661,0.025803) (0.022449,0.023918) (0.019449,0.022056) (0.016662,0.02022) (0.014087,0.018413) (0.011727,0.01664) (0.0095815,0.014902) (0.0076509,0.013204) (0.0059358,0.011546) (0.0044365,0.009932) (0.0031534,0.0083637) (0.0020865,0.006843) (0.0012362,0.0053715) (0.00060239,0.0039505) (0.00018514,0.0025811) (-1.5618e-05,0.0012641) (0,0) (0.00023493,-0.0012421) (0.00069203,-0.0024934) (0.0013709,-0.0037534) (0.0022709,-0.0050215) (0.0033916,-0.0062971) (0.0047319,-0.0075791) (0.0062911,-0.0088664) (0.0080682,-0.010158) (0.010062,-0.011452) (0.012271,-0.012747) (0.014694,-0.014041) (0.017329,-0.015333) (0.020176,-0.016621) (0.023232,-0.017901) (0.026494,-0.019173) (0.029963,-0.020434) (0.033635,-0.021681) (0.037508,-0.022913) (0.041581,-0.024126) (0.04585,-0.025318) (0.050314,-0.026487) (0.054971,-0.027631) (0.059817,-0.028745)};
\addplot[cblue,very thick] coordinates {(0.058521,0.045046) (0.053639,0.043279) (0.04896,0.041488) (0.044485,0.039677) (0.040216,0.037846) (0.036156,0.035998) (0.032304,0.034135) (0.028664,0.032259) (0.025236,0.030372) (0.022021,0.028476) (0.019021,0.026573) (0.016237,0.024665) (0.01367,0.022752) (0.011321,0.020838) (0.0091904,0.018923) (0.0072788,0.017008) (0.005587,0.015096) (0.0041153,0.013188) (0.0028643,0.011284) (0.0018341,0.0093849) (0.001025,0.0074925) (0.00043712,0.0056073) (7.0386e-05,0.0037299) (-7.5258e-05,0.0018606) (0,0) (0.00029457,-0.0018387) (0.00080678,-0.0036422) (0.0015362,-0.0054103) (0.0024821,-0.0071426) (0.003644,-0.0088389) (0.005021,-0.010499) (0.0066124,-0.012122) (0.008417,-0.013708) (0.010434,-0.015257) (0.012662,-0.016767) (0.0151,-0.018239) (0.017746,-0.019672) (0.0206,-0.021066) (0.023659,-0.022419) (0.026922,-0.023732) (0.030388,-0.025003) (0.034054,-0.026233) (0.037919,-0.027421) (0.041981,-0.028566) (0.046238,-0.029668) (0.050688,-0.030726) (0.055328,-0.031739)};
\addplot[cmuted,very thick] coordinates {(0.0582,0.048848) (0.053299,0.047242) (0.048603,0.045597) (0.044112,0.043915) (0.039828,0.042195) (0.035755,0.040438) (0.031893,0.038643) (0.028244,0.036811) (0.02481,0.034942) (0.021593,0.033035) (0.018594,0.031091) (0.015813,0.02911) (0.013254,0.027091) (0.010915,0.025036) (0.0087994,0.022943) (0.0069068,0.020813) (0.0052381,0.018647) (0.0037941,0.016443) (0.0025752,0.014203) (0.0015817,0.011927) (0.00081385,0.0096136) (0.00027184,0.0072642) (-4.4368e-05,0.0048787) (-0.0001349,0.0024572) (0,0) (0.00035421,-0.0024353) (0.00092154,-0.004791) (0.0017014,-0.0070671) (0.0026933,-0.0092636) (0.0038964,-0.011381) (0.0053101,-0.013419) (0.0069336,-0.015378) (0.0087658,-0.017259) (0.010806,-0.019062) (0.013053,-0.020788) (0.015506,-0.022437) (0.018163,-0.024011) (0.021024,-0.025511) (0.024087,-0.026937) (0.02735,-0.02829) (0.030813,-0.029573) (0.034474,-0.030785) (0.038331,-0.031929) (0.042382,-0.033006) (0.046626,-0.034017) (0.051061,-0.034964) (0.055686,-0.035848)};
\end{groupplot}
\end{tikzpicture}

\caption{NACA 2412 modifié pour $I=2$, 4, 6, 8 (calculé par \code{naca4\_modified}). L'épaisseur maximale et sa position ne changent pas ; seul le nez, et la façon dont l'épaisseur y est atteinte, changent.}\label{fig:v2-formes}
\end{figure}

\begin{codebox}[ : \code{aero/naca4.py}, V2]
\begin{lstlisting}
def _modified_coeffs(I, T):
    d1 = (2.24 - 5.42*T + 12.3*T**2) / (10*(1 - 0.878*T))          # pente au bord de fuite
    u = 1 - T
    d2, d3 = np.linalg.solve([[u**2, u**3], [2*u, 3*u**2]], [0.1 - d1*u, -d1])
    ypp = 2*d2 + 6*d3*u                                           # courbure de l'arc arrière en T
    a0 = 0.296904 * I / 6
    M = [[T, T**2, T**3], [1, 2*T, 3*T**2], [0, 2, 6*T]]
    rhs = [0.1 - a0*np.sqrt(T), -0.5*a0/np.sqrt(T), ypp + 0.25*a0*T**-1.5]
    a1, a2, a3 = np.linalg.solve(M, rhs)
\end{lstlisting}
\end{codebox}

\section{XFOIL comme simulateur de référence}
\paragraph{Pourquoi.} Le test V5 (§\ref{sec:v5}) a montré que le modèle rapide classe bien les profils (corrélation de rang 0,91) mais surestime $C_D$ de 27~\% et déplace l'optimum. Pour un livrable, il faut le simulateur le plus juste disponible ; XFOIL coûte environ une seconde par point, et la grille V2 (1\,024 profils) se calcule en 6~minutes sur 10 cœurs. Le modèle rapide reste utile pour développer et tester, là où sa vitesse compte.

\paragraph{Le point de croisière, à portance imposée.} On impose $C_L=0{,}4$ (commande \texttt{CL} d'XFOIL) : on compare ainsi les profils à mission égale, et XFOIL renvoie l'incidence de croisière, qui est aussi l'attitude vue par le radar. Le $C_D$ correspondant est le premier objectif.

\paragraph{La contrainte de manœuvre.} Un avion doit pouvoir manœuvrer : on exige une portance maximale $C_{L,\max}\ge1{,}2$, estimée par une rampe d'incidence de 4 à $14^\circ$. Un nez aigu décroche tôt : la couche limite décolle juste après le pic d'aspiration, très intense sur un petit rayon. Dans XFOIL, cela se voit de deux façons : une portance qui plafonne puis chute, ou une non-convergence au-delà d'une certaine incidence. Dans les deux cas, $C_{L,\max}$ est pris comme la plus grande portance convergée.

\begin{figure}[htbp]
\centering
% figure générée par tools/make_figures_v2.py : v2_xfoil
\begin{tikzpicture}
\begin{groupplot}[group style={group size=2 by 1,horizontal sep=1.6cm},width=0.48\linewidth,height=5.4cm,
  xlabel={épaisseur $t$ (\%)},title style={font=\small},tick label style={font=\scriptsize},legend style={draw=none,font=\scriptsize}]
\nextgroupplot[title={$C_D$ de croisière ($C_L=0{,}4$), médiane sur $m,p$},ylabel={counts},legend pos=north west]
\addplot[corange,very thick,mark=*,mark size=1.2pt] coordinates {(8,48) (10,50.8) (12,54.1) (14,59.6) (16,61.6) (18,65.3) (20,70.3)};\addlegendentry{$I=2$}
\addplot[cgreen,very thick,mark=*,mark size=1.2pt] coordinates {(8,53.5) (10,52.9) (12,56.1) (14,58) (16,60.6) (18,64.7) (20,69.3)};\addlegendentry{$I=4$}
\addplot[cblue,very thick,mark=*,mark size=1.2pt] coordinates {(8,56) (10,55.9) (12,57.7) (14,59.7) (16,62.45) (18,65.3) (20,70.15)};\addlegendentry{$I=6$}
\addplot[cmuted,very thick,mark=*,mark size=1.2pt] coordinates {(8,55.95) (10,59.2) (12,61.25) (14,63.3) (16,67.5) (18,70.05) (20,76.35)};\addlegendentry{$I=8$}
\nextgroupplot[title={$C_{L,\max}$ (rampe 4--14°), médiane sur $m,p$},ymin=0.6]
\addplot[corange,very thick,mark=*,mark size=1.2pt] coordinates {(8,0.981) (10,1.0245) (12,1.0112) (14,1.1883) (16,1.328) (18,1.3478) (20,1.2788)};
\addplot[cgreen,very thick,mark=*,mark size=1.2pt] coordinates {(8,1.2839) (10,1.5288) (12,1.5613) (14,1.5372) (16,1.5132) (18,1.4736) (20,1.5252)};
\addplot[cblue,very thick,mark=*,mark size=1.2pt] coordinates {(8,1.6178) (10,1.6147) (12,1.609) (14,1.5857) (16,1.5264) (18,1.5132) (20,1.5011)};
\addplot[cmuted,very thick,mark=*,mark size=1.2pt] coordinates {(8,1.6741) (10,1.6447) (12,1.6326) (14,1.6006) (16,1.5582) (18,1.4953) (20,1.4599)};
\addplot[cink,dashed] coordinates {(8,1.2) (20,1.2)};
\node[font=\scriptsize,anchor=south east] at (axis cs:20,1.2) {exigence de manœuvre};
\end{groupplot}
\end{tikzpicture}

\caption{Effet du rayon de bord d'attaque sur les 1\,024 profils V2 (XFOIL, $\Re=3\cdot10^6$). Un nez aigu ($I=2$) réduit un peu la traînée de croisière, mais fait chuter la portance maximale sous l'exigence de manœuvre.}\label{fig:v2-xfoil}
\end{figure}

\paragraph{Ce que montre la grille.}
\begin{itemize}
\item À $t=12$~\%, en médiane sur la cambrure, passer de $I=8$ à $I=2$ \emph{réduit} la traînée de croisière de 61 à 54~counts : un nez plus fin limite la survitesse et le gradient adverse qui suit.
\item Mais $C_{L,\max}$ tombe de 1,63 à 1,01. Parmi les profils assez épais, 94 profils à $I=2$ ne tiennent pas l'exigence de manœuvre, contre 22 à $I=4$ et moins de 10 au-delà.
\item XFOIL ne converge pas en croisière sur 104 profils, dont 70 à $I=2$ : ils sont déclarés infaisables. Au total, 702 profils sur 1\,024 sont faisables.
\end{itemize}

\begin{projet}
C'est exactement le mécanisme qui manquait en V1. La SER veut un nez aigu (cours SER, chapitre V2), la croisière aussi, mais la manœuvre l'interdit. Le compromis traînée/SER se joue désormais sur l'épaisseur, la cambrure \emph{et} le rayon de bord d'attaque, et le front de Pareto n'est plus confiné à $t=t_{\min}$ (cours d'optimisation, chapitre V2).
\end{projet}

\begin{piege}
La non-convergence d'XFOIL n'est pas une preuve de décrochage : elle peut venir d'un bulbe laminaire mal résolu ou d'un manque d'itérations. Considérer « non convergé » comme « décroché » est une hypothèse prudente mais grossière. Un critère plus propre serait une polaire complète avec détection du maximum de $C_L$, ou un calcul RANS sur quelques profils de contrôle.
\end{piege}

\exercices
\begin{exercice}
Pour $t=12$~\%, calculer le rayon de bord d'attaque pour $I=2$, 4, 6 et 8. De combien de décibels la SER de nez (loi $\pi r_{BA}$) diminue-t-elle entre $I=8$ et $I=2$ ?
\end{exercice}

\begin{pistes}
\begin{itemize}
\item Libérer aussi la position de l'épaisseur maximale $T$ (cinquième variable), qui agit sur la position de la transition, donc sur la traînée.
\item Remplacer « non convergé = décroché » par une détection propre du maximum de $C_L$ sur une polaire fine, et valider sur quelques profils par RANS.
\item Calibrer le modèle rapide sur XFOIL (seuil de Michel, décambrage visqueux) pour disposer d'un simulateur basse fidélité fiable, utile à une approche multi-fidélité.
\end{itemize}
\end{pistes}
